\Needspace{18\baselineskip}
\begin{center}\textit{Some special theorems.}\end{center}
\SixthArticle{53}
All the relations we have derived so far are most general in that the elements $\alpha,\beta,\gamma$ are subject to no conditions. We have, however, also found several others that presuppose special conditions between the elements $\alpha,\beta,\gamma$: undoubtedly many more still lie hidden, and perhaps even those we shall give here may in future be derived from higher principles.
\par\SourceLeaf{233}
First, in equation 80, let us set $x=\frac{4y}{(1+y)^2}$; then
\[
\mathrm dx=\mathrm dy\mathbin{.}\frac{4(1-y)}{(1+y)^3}
\]
and therefore
\[
\begin{aligned}
\frac{\mathrm dP}{\mathrm dx}&=\frac{\mathrm dP}{\mathrm dy}\mathbin{.}\frac{(1+y)^3}{4(1-y)}\\
\frac{\mathrm{dd}P}{\mathrm dx}&=\mathrm d\frac{\mathrm dP}{\mathrm dy}\mathbin{.}\frac{(1+y)^3}{4(1-y)}\\&\qquad+\frac{\mathrm dP}{\mathrm dy}\mathbin{.}\frac{(2-y)(1+y)^2}{2(1-y)^2}\mathrm dy\\
\frac{\mathrm{dd}P}{\mathrm dx^2}&=\frac{\mathrm{dd}P}{\mathrm dy^2}\mathbin{.}\frac{(1+y)^6}{16(1-y)^2}\\&\qquad+\frac{(2-y)(1+y)^5}{8(1-y)^3}\mathbin{.}\frac{\mathrm dP}{\mathrm dy}
\end{aligned}
\]
That equation consequently becomes
\[
\begin{aligned}
0={}&\alpha\beta P\\
&-(\gamma(1+y)^2-4(\alpha+\beta+1)y)\frac{1+y}{4(1-y)}\mathbin{.}\frac{\mathrm dP}{\mathrm dy}\\
&-\frac{\mathrm{dd}P}{\mathrm dy^2}\mathbin{.}\frac{y(1+y)^2}4
-\frac{y(2-y)(1+y)}{2(1-y)}\mathbin{.}\frac{\mathrm dP}{\mathrm dy}
\end{aligned}
\]
or
\[
\begin{aligned}
&0=4\alpha\beta(1-y)P\\
&-(\gamma(1+y)^2-4(\alpha+\beta+1)y-2y(2-y)(1+y))\frac{\mathrm dP}{\mathrm dy}\\
&-(y-yy)(1+y)^2\frac{\mathrm{dd}P}{\mathrm dy^2}\\
&0=4\alpha\beta(1-y)P\\
&-(1+y)(\gamma-(4\alpha+4\beta-2\gamma)y+(\gamma-2)yy)\frac{\mathrm dP}{\mathrm dy}\\
&-(1+y)^2(y-yy)\frac{\mathrm{dd}P}{\mathrm dy^2}
\end{aligned}
\]
Setting $P=(1+y)^{2\alpha}Q$ gives
\[
\begin{aligned}
\mathrm I.\quad 0={}&2\alpha(2\beta-\gamma+(2\alpha+1-\gamma)y)Q\\
&-(\gamma-(4\beta-2\gamma)y+(\gamma-4\alpha-2)yy)\frac{\mathrm dQ}{\mathrm dy}\\
&-(y-yy)(1+y)\frac{\mathrm{dd}Q}{\mathrm dy^2}
\end{aligned}
\]
Now, on supposing $\beta=\alpha+\frac12$, this equation takes the following form:
\[
\begin{aligned}
0={}&2\alpha(2\alpha+1-\gamma)Q\\
&-(\gamma-(4\alpha+2-\gamma)y)\frac{\mathrm dQ}{\mathrm dy}\\
&-(y-yy)\frac{\mathrm{dd}Q}{\mathrm dy^2}
\end{aligned}
\]
whose solution is
\[
Q=F(2\alpha,2\alpha+1-\gamma,\gamma,y)
\]
\SourceLeaf{234}
so that the result is
\[
\begin{aligned}&[100]\quad(1+y)^{2\alpha}F(2\alpha,2\alpha+1-\gamma,\gamma,y)\\&=F(\alpha,\alpha+\tfrac12,\gamma,\tfrac{4y}{(1+y)^2})\end{aligned}
\]
\SixthArticle{54}
If instead of the relation $\beta=\alpha+\frac12$ we adopt $\gamma=2\beta$, equation I of the preceding article becomes
\[
\begin{aligned}
0={}&2\alpha(2\alpha+1-2\beta)yQ\\
&-(2\beta-(4\alpha+2-2\beta)yy)\frac{\mathrm dQ}{\mathrm dy}\\
&-y(1-yy)\frac{\mathrm{dd}Q}{\mathrm dy^2}
\end{aligned}
\]
Now setting $yy=z$ gives
\[
\begin{aligned}
\frac{\mathrm dQ}{\mathrm dy}&=2y\frac{\mathrm dQ}{\mathrm dz}\\
\frac{\mathrm{dd}Q}{\mathrm dy^2}&=4yy\frac{\mathrm{dd}Q}{\mathrm dz^2}+\frac{2\mathrm dQ}{\mathrm dz}
\end{aligned}
\]
and therefore
\[
\begin{aligned}
0={}&\alpha(\alpha+\tfrac12-\beta)Q\\
&-(\beta+\tfrac12-(2\alpha+\tfrac32-\beta)z)\frac{\mathrm dQ}{\mathrm dz}\\
&-(z-zz)\frac{\mathrm{dd}Q}{\mathrm dz^2}
\end{aligned}
\]
whose solution is
\[
Q=F(\alpha,\alpha+\tfrac12-\beta,\beta+\tfrac12,z)
\]
so we have
\[
\begin{aligned}&[101]\quad(1+y)^{2\alpha}F(\alpha,\alpha+\tfrac12-\beta,\beta+\tfrac12,yy)\\&=F(\alpha,\beta,2\beta,\tfrac{4y}{(1+y)^2})\end{aligned}
\]
\SixthArticle{55}
Secondly, let us set $x=4y-4yy$; then
\[
\begin{aligned}
\mathrm dx&=4\mathrm dy(1-2y)\\
\frac{\mathrm dP}{\mathrm dx}&=\frac{\mathrm dP}{\mathrm dy}\mathbin{.}\frac1{4(1-2y)}\\
\frac{\mathrm{dd}P}{\mathrm dx^2}&=\frac{\mathrm{dd}P}{\mathrm dy^2}\mathbin{.}\frac1{16(1-2y)^2}
+\frac{\mathrm dP}{\mathrm dy}\mathbin{.}\frac1{8(1-2y)^3}
\end{aligned}
\]
and equation 80 consequently becomes
\SourceLeaf{235}
\[
\begin{aligned}
&0=4\alpha\beta P\\
&-(\gamma-(4\alpha+4\beta+2)y+(4\alpha+4\beta+2)yy)\frac1{(1-2y)}\mathbin{.}\frac{\mathrm dP}{\mathrm dy}\\
&-(y-yy)\frac{\mathrm{dd}P}{\mathrm dy^2}
\end{aligned}
\]
To remove the fraction in the second term, $\gamma=\alpha+\beta+\frac12$ must be set, giving
\[
\begin{aligned}
0={}&4\alpha\beta P\\
&-(\alpha+\beta+\tfrac12-(2\alpha+2\beta+1)y)\frac{\mathrm dP}{\mathrm dy}\\
&-(y-yy)\frac{\mathrm{dd}P}{\mathrm dy^2}
\end{aligned}
\]
whose solution is
\[
P=F(2\alpha,2\beta,\alpha+\beta+\tfrac12,y)
\]
so we have
\[
\begin{aligned}&[102]\quad F(\alpha,\beta,\alpha+\beta+\tfrac12,4y-4yy)\\&=F(2\alpha,2\beta,\alpha+\beta+\tfrac12,y)\end{aligned}
\]
If in this equation we were to change $y$ into $1-y$, the result would be
\[
F(\alpha,\beta,\alpha+\beta+\tfrac12,4y-4yy)=F(2\alpha,2\beta,\alpha+\beta+\tfrac12,1-y)
\]
from which the paradox seems to follow:
\[
F(2\alpha,2\beta,\alpha+\beta+\tfrac12,y)=F(2\alpha,2\beta,\alpha+\beta+\tfrac12,1-y)
\]
This equation is certainly false. To resolve this, we must remember to distinguish carefully between the two meanings of the symbol $F$, according as it represents \textit{either} the function whose nature is expressed by differential equation 80, \textit{or} merely the sum of the infinite series. The latter always gives a fully determined quantity as long as the fourth element lies between $-1$ and $+1$; but you must take care not to exceed these bounds, since otherwise no meaning at all remains. The former meaning, however, represents a general function, which always changes according to the law of continuity when the fourth element changes in a continuous flow, whether you assign it real or imaginary values, provided that you always avoid the values 0 and 1. It follows that, in the latter sense, the function can acquire unequal values for equal values of the fourth element, after passing, or rather returning, through imaginary quantities; the value represented by the \textit{series} $F$
\SourceLeaf{236}
is only one of these. There is therefore no contradiction at all in \textit{some} value of the function $F(\alpha,\beta,\alpha+\beta+\frac12,4y-4yy)$ being equal to $F(2\alpha,2\beta,\alpha+\beta+\frac12,y)$, while \textit{another} value becomes $=F(2\alpha,2\beta,\alpha+\beta+\frac12,1-y)$; and it would be just as absurd to infer from this that these values are equal as to infer from $\mathrm{Arc.\ sin}\frac12=30^\circ$, $\mathrm{Arc.\ sin}\frac12=150^\circ$ that $30^\circ=150^\circ$. --- But if we take the symbol $F$ in the less general meaning, namely as representing only the sum of the series $F$, the arguments by which we derived equation 102 necessarily suppose that $y$ increases from the value 0 only until $x=1$ is reached, that is, as far as $y=\frac12$. At this very point, however, the \textit{continuity} of the series $P=F(\alpha,\beta,\alpha+\beta+\frac12,4y-4yy)$ is interrupted, since $\frac{\mathrm dP}{\mathrm dy}$ clearly jumps suddenly from a positive (finite) value to a negative one. Thus, in this meaning, equation 102 cannot be extended beyond the bounds from $y=\frac12-\BareRoot{\frac12}$ to $y=\frac12$. If you prefer, you can also express the same equation as follows:
\[
\begin{aligned}&[103]\quad F(\alpha,\beta,\alpha+\beta+\tfrac12,x)\\&=F(2\alpha,2\beta,\alpha+\beta+\tfrac12,\tfrac{1-\sqrt{1-x}}2)\end{aligned}
\]
or as follows:
\[
\begin{aligned}&[104]\quad F(\alpha,\beta,\alpha+\beta+\tfrac12,1-x)\\&=F(2\alpha,2\beta,\alpha+\beta+\tfrac12,\tfrac{1-\BareRoot{x}}2)\end{aligned}
\]
which yields the following corollary (formula 48):
\[
\begin{aligned}
{}[105]\quad F(2\alpha,2\beta,\alpha+\beta+\tfrac12,\tfrac12)&=\frac{\Pi(\alpha+\beta-\tfrac12)\Pi(-\tfrac12)}{\Pi(\alpha-\tfrac12)\Pi(\beta-\tfrac12)}\\
&=\frac{\Pi(\alpha+\beta-\tfrac12)\BareRoot{\pi}}{\Pi(\alpha-\tfrac12)\Pi(\beta-\tfrac12)}
\end{aligned}
\]
\SixthArticle{56}
Applying formula 87 to equation 104 gives
\[
\begin{aligned}&[106]\quad F(2\alpha,2\beta,\alpha+\beta+\tfrac12,\tfrac{1-\BareRoot{x}}2)\\&=AF(\alpha,\beta,\tfrac12,x)+B\BareRoot{x}\mathbin{.}F(\alpha+\tfrac12,\beta+\tfrac12,\tfrac32,x)\end{aligned}
\]
which shows that the series
\[
F(2\alpha,2\beta,\alpha+\beta+\tfrac12,\tfrac{1-t}2)
\]
can be represented by the series
\[
\begin{aligned}&A+Bt+A\frac{\alpha\mathbin{.}\beta}{1\mathbin{.}\tfrac12}\mathbin{.}tt\\&+B\frac{\alpha+\tfrac12\mathbin{.}\beta+\tfrac12}{1\mathbin{.}\tfrac32}t^3\\&+A\frac{\alpha\mathbin{.}\alpha+1\mathbin{.}\beta\mathbin{.}\beta+1}{1\mathbin{.}2\mathbin{.}\tfrac12\mathbin{.}\tfrac32}t^4+\text{etc.}\end{aligned}
\]
setting, for brevity,
\SourceLeaf{237}
\[
\begin{aligned}&A=\frac{\Pi(\alpha+\beta-\tfrac12)\Pi(-\tfrac12)}{\Pi(\alpha-\tfrac12)\Pi(\beta-\tfrac12)},\\&
B=\frac{\Pi(\alpha+\beta-\tfrac12)\Pi(-\tfrac32)}{\Pi(\alpha-1)\Pi(\beta-1)}\end{aligned}
\]
From this it may be concluded that
\[
\begin{aligned}&[107]\quad F(2\alpha,2\beta,\alpha+\beta+\tfrac12,\tfrac{1+\BareRoot{x}}2)\\&=AF(\alpha,\beta,\tfrac12,x)-B\BareRoot{x}\mathbin{.}F(\alpha+\tfrac12,\beta+\tfrac12,\tfrac32,x)\end{aligned}
\]
If anyone should find this conclusion insufficiently justified (although it would not be difficult to put it beyond all doubt), we could reach the same equation in the following way. Equation 87 gives
\[
\begin{aligned}
&F(2\alpha,2\beta,\alpha+\beta+\tfrac12,\tfrac{1+\BareRoot{x}}2)\\&=CF(2\alpha,2\beta,\alpha+\beta+\tfrac12,\tfrac{1-\BareRoot{x}}2)\\
&\quad+D\left(\frac{1-x}4\right)^{\frac12-\alpha-\beta}F(1-2\alpha,1-2\beta,\tfrac32-\alpha-\beta,\tfrac{1-\BareRoot{x}}2)
\end{aligned}
\]
setting, for brevity,
\[
\begin{aligned}&C=\frac{\Pi(\alpha+\beta-\tfrac12)\Pi(-\tfrac12-\alpha-\beta)}{\Pi(\alpha-\beta-\tfrac12)\Pi(\beta-\alpha-\tfrac12)},\\&
D=\frac{\Pi(\alpha+\beta-\tfrac12)\Pi(\alpha+\beta-\tfrac32)}{\Pi(2\alpha-1)\Pi(2\beta-1)}\end{aligned}
\]
But equation 104 easily gives
\[
\begin{aligned}
&F(1-2\alpha,1-2\beta,\tfrac32-\alpha-\beta,\tfrac{1-\BareRoot{x}}2)\\&=F(\tfrac12-\alpha,\tfrac12-\beta,\tfrac32-\alpha-\beta,1-x)\\
&=EF(\tfrac12-\alpha,\tfrac12-\beta,\tfrac12,x)\\&\qquad+G\BareRoot{x}\mathbin{.}F(1-\alpha,1-\beta,\tfrac32,x)
\end{aligned}
\]
setting, for brevity,
\[
\begin{aligned}&E=\frac{\Pi(\tfrac12-\alpha-\beta)\Pi(-\tfrac12)}{\Pi(-\alpha)\Pi(-\beta)},\\&
G=\frac{\Pi(\tfrac12-\alpha-\beta)\Pi(-\tfrac32)}{\Pi(-\tfrac12-\alpha)\Pi(-\tfrac12-\beta)}\end{aligned}
\]
From this, equation 82 again gives
\[
\begin{aligned}
&F(1-2\alpha,1-2\beta,\tfrac32-\alpha-\beta,\tfrac{1-\BareRoot{x}}2)\\
&\quad=E(1-x)^{\alpha+\beta-\frac12}F(\alpha,\beta,\tfrac12,x)\\&\qquad+G\BareRoot{x}\mathbin{.}(1-x)^{\alpha+\beta-\frac12}F(\alpha+\tfrac12,\beta+\tfrac12,\tfrac32,x)
\end{aligned}
\]
On substituting these, and setting
\[
AC+DE2^{2\alpha+2\beta-1}=M,\quad BC+DG2^{2\alpha+2\beta-1}=N
\]
the result is
\[
\begin{aligned}&F(2\alpha,2\beta,\alpha+\beta+\tfrac12,\tfrac{1+\BareRoot{x}}2)\\&=MF(\alpha,\beta,\tfrac12,x)+N\BareRoot{x}\mathbin{.}F(\alpha+\tfrac12,\beta+\tfrac12,\tfrac32,x)\end{aligned}
\]
whose \textit{form} agrees with equation 107. Now we could indeed derive from the nature of the function $\Pi$ alone that $M=A$, $N=-B$, since it is easily shown by equations 55, 56 that
\[
\begin{aligned}&C=\frac{\cos(\alpha-\beta)\pi}{\cos(\alpha+\beta)\pi},\\&
\frac{DE2^{2\alpha+2\beta-1}}A=-\frac{2\sin\alpha\pi\sin\beta\pi}{\cos(\alpha+\beta)\pi},\\&
\frac{DG2^{2\alpha+2\beta-1}}B=-\frac{2\cos\alpha\pi\cos\beta\pi}{\cos(\alpha+\beta)\pi},\end{aligned}
\]
\SourceLeaf{238}
but even this work is unnecessary. For setting $x=0$ clearly requires
\[
M=F(2\alpha,2\beta,\alpha+\beta+\tfrac12,\tfrac12)=A
\]
while differentiating that equation gives
\[
\begin{aligned}
&x^{-\frac12}\frac{\alpha\beta}{\alpha+\beta+\tfrac12}F(2\alpha+1,2\beta+1,\alpha+\beta+\tfrac32,\tfrac{1+\BareRoot{x}}2)\\
&\quad=2\alpha\beta MF(\alpha+1,\beta+1,\tfrac32,x)\\
&\qquad+\tfrac23(\alpha+\tfrac12)(\beta+\tfrac12)N\BareRoot{x}\mathbin{.}F(\alpha+\tfrac32,\beta+\tfrac32,\tfrac52,x)\\&\qquad+\tfrac12Nx^{-\frac12}F(\alpha+\tfrac12,\beta+\tfrac12,\tfrac32,x)
\end{aligned}
\]
from which, on setting $x=0$, follows
\[
\begin{aligned}
N&=\frac{2\alpha\beta}{\alpha+\beta+\tfrac12}F(2\alpha+1,2\beta+1,\alpha+\beta+\tfrac32,\tfrac12)\\
&=\frac{2\alpha\beta}{\alpha+\beta+\tfrac12}\mathbin{.}\frac{\Pi(\alpha+\beta+\tfrac12)\Pi(-\tfrac12)}{\Pi\alpha\Pi\beta}\\
&=\frac{\Pi(\alpha+\beta-\tfrac12)\Pi(-\tfrac32)}{\Pi(\alpha-1)\Pi(\beta-1)}=-B
\end{aligned}
\]
